\documentclass{article}
\usepackage[UTF8]{ctex}
\usepackage{amsmath,amssymb,amsfonts}
\usepackage{enumitem}
\usepackage{geometry}
\geometry{a4paper, margin=1in}
% === 错题本排版优化（不参与题目标识） ===
\setlength{\parindent}{0pt}
\setlength{\parskip}{0.35em plus 0.10em minus 0.08em}
\setlength{\emergencystretch}{3em}
\setlist{itemsep=0.25em, topsep=0.40em, parsep=0pt}
\allowdisplaybreaks[2]
% === 错题本排版优化结束 ===
\begin{document}
\section*{山东大学2026年数学分析真题}
\section*{123．山东大学2026年数学分析试卷}

\subsection{题目1}
1 ．求极限 \(\lim _{n \rightarrow \infty} n\left(e-\left(1+\frac{1}{n}\right)^{n}\right)\) ．

\subsection{题目2}
2．设 \(f(x)=x^{\alpha}, x \in[1,+\infty)\) ，求 \(\alpha\) 的取值范围使得 \(f(x)\) 在 \([1,+\infty)\) 上一致连续．

\subsection{题目3}
3．已知定义在 \(D=\left\{(x, y) \mid x^{2}+y^{2} \leq 1\right\}\) 上的二阶连续可微函数 \(f(x, y)\) 满足 \(\frac{\partial^{2} f}{\partial x^{2}}+\frac{\partial^{2} f}{\partial y^{2}}=x^{2}+y^{2}\) ．计算

\[
\iint_{D}\left(\frac{x}{\sqrt{x^{2}+y^{2}}} \frac{\partial f}{\partial x}+\frac{y}{\sqrt{x^{2}+y^{2}}} \frac{\partial f}{\partial y}\right) \mathrm{d} x \mathrm{~d} y .
\]

\subsection{题目4}
4．已知 \(k>0\) ，若 \(\lim _{n \rightarrow \infty} \frac{b_{1}+b_{2}+\cdots+b_{n}}{n^{k}}=b\) ，求证： \(\lim _{n \rightarrow \infty} \frac{b_{n}}{n^{k}}=0\) ．

\subsection{题目5}
5．已知 \(f(x) \in C[0,2 \pi]\) ，证明： \(\lim _{n \rightarrow \infty} \int_{0}^{2 \pi} f(x)|\sin n x| \mathrm{d} x=\frac{2}{\pi} \int_{0}^{2 \pi} f(x) \mathrm{d} x\) ．

\subsection{题目6}
6．设 \(f(x)=\sum_{n=1}^{\infty} \frac{x^{n}}{n^{2} \ln (1+n)}\) ，讨论 \(f(x)\) 在 \([-1,1]\) 上的连续性，以及在 \(x=-1\) 和 \(x=1\) 时的可微性．

\subsection{题目7}
7．设 \(f(x)=x^{3}\left(e^{\frac{1}{x}}+e^{-\frac{1}{x}}-c\right)\) 满足 \(\lim _{x \rightarrow \infty}[f(x)-(a x+b)]=0\) ，求 \(a, b, c\) ，并说明其几何意义。

\subsection{题目8}
8．设 \(f(t)\) 在 \([0,+\infty)\) 上连续，且

\[
f(t)=\iint_{x^{2}+y^{2} \leq t^{2}}\left(x^{2}+y^{2}\right) f\left(\sqrt{x^{2}+y^{2}}\right) \mathrm{d} x \mathrm{~d} y+t^{4}, t>0 .
\]

求 \(f(t)\) ．

\subsection{题目9}
9．已知球面 \(\Gamma: x^{2}+y^{2}+z^{2}=a^{2}(a>0)\) ，球面 \(\sum\) 的球心在 \(\Gamma\) 上，半径为 \(R\) ，当 \(R\) 为何值时，球面 \(\sum\)位于 \(\Gamma\) 内的面积取最大值？并求出最大面积。
\end{document}
